\scanpage{526}
\mspage{437}
% B01
fibrée, resp. une sous-catégorie cofibrée, sur $C$ (le foncteur $E(\pi)\unsure{\relbar|\to}E$ étant donc cartésien resp cocartésien). À $C$-équivalence près, $E$ se reconstitue alors \unsure{:|à} partir de $E(\pi)$, en passant aux $\mathrm{Ind}_\pi(E_{s\pi})$, \deleted{\unsure{et on obtient}} qui forment un système fibré resp. cofibré sur $C$ par les foncteurs $\mathrm{Ind}_\pi(\lambda^\ast_\pi)$ resp $\mathrm{Ind}_\pi(\lambda_*^\pi)$ ....

% B02
% slanted
\margin{Ici on utilise seulement $\pi\geq\mathrm{card\,red}\,C$.}

% B03
Prouver que $E$ satisfait à $L_\pi$ et que $f:E\to C$ est $\pi$-accessible. Comme $C$ est karoubienne, et $\pi\geq\pi_1\geq\mathrm{Card}\,C$, $C$ satisfait à $L_\pi$ et les \deleted{\unsure{systèmes}} \unsure{``ind-objet|``ind-objets|``$-1$-déplt} \unsure{grds|pet} devant $\pi$ de $C$ sont ess. constants \inserted{(4.8.13)}.

% B04
Soit alors $I$ \unsure{grd|filt} devant $\pi$, et $I\longrightarrow E$, $i\mapsto x_i$, un foncteur, il faut construire une $\underline{\lim}$ des $x_i$ dans $E$, en utilisant la $\varinjlim$ déjà connue $s$ des $s_i=f(x_i)$. \deleted{\unsure{Il}} On va prendre \deleted{\unsure{son}} $x$ dans $E_s$,
% B05
\diagram{\begin{tikzcd}[ampersand replacement=\&]x_i\arrow[r,"\alpha_i"]\& x\ {?}\end{tikzcd}}

% B06
Prenons d'abord le cas où $f$ est cofibrant. On voit que les \deleted{\unsure{$x_i$}} $\beta_{i*}(x_i)$ forment un système inductif dans $E_s$, puisque pour
\diagram{\begin{tikzcd}[ampersand replacement=\&,column sep=small,row sep=small]s_i\arrow[rr,"\beta_i"]\arrow[dr,"\beta_{ji}"']\&{}\&s\\ \&s_j\arrow[ur,"\beta_j"']\&{}\end{tikzcd}}
$i\overset{\deleted{\unsure{\alpha_{ji}}}}{\unsure{\leq|\longrightarrow}}j$, on a
% B07
\[ \beta_{i*}(\unsure{x_j|x_i})=\beta_{j*}\underset{\longmapsto x_j}{\underset{\rule{.35pt}{.4ex}\rule{4.9em}{.35pt}\rule{.35pt}{.4ex}}{\beta_{ji*}(x_{\unsure{j|i}})}}\longrightarrow\beta_{j*}(x_j). \]

% B08
On \deleted{peut \unsure{y}} pose

% archive sheet number: 525
